\section{从题目求解到数学研究}

AI for Math 可以分成不同层级：做题、形式化证明、证明补全、猜想生成、研究协作。越往后，越接近真正数学研究，也越难评估。竞赛题有答案，Lean 证明有 verifier，研究问题则需要判断“值不值得做”。

\begin{figure}[H]
\centering
\includegraphics[width=0.96\textwidth]{ai-for-math-roadmap.png}
\caption{从数学题到数学研究的 AI for Math 路线图。}
\end{figure}

\begin{knowledgebox}{读图：为什么越往右越难}
题目求解通常有答案；形式化证明有 proof assistant 检查；证明补全有局部目标；猜想生成和研究协作则需要审美、方向判断和长期理论结构。越接近研究，越不能只靠 benchmark。
\end{knowledgebox}

\subsection{数学家的直觉}

访谈后半段讨论“数学家的直觉”。直觉不是玄学，而是长期训练后对结构、反例、证明路线和审美的压缩判断。AI 可以通过大量数据和反馈学习模式，但是否能形成类似数学家直觉，取决于它是否能把形式化反馈、非形式推理和研究目标整合起来。

\begin{knowledgebox}{直觉可以被训练，但不等于数据拟合}
数学直觉来自问题选择、失败经验、结构类比和审美判断。AI 可以从证明库、论文、交互反馈和数学家修正中学习，但需要长期闭环，而不是一次性监督数据。
\end{knowledgebox}

\subsection{登月要么成功，要么失败}

访谈中“登月要么成功，要么失败”的表达，体现了 AI for Math 创业的高风险性质。数学研究不像普通应用，可以小步验证商业指标；它要么真的推动证明和发现能力，要么只是漂亮 demo。Axiom 这类 Neo Lab 的难度就在于：要同时做科学、产品、团队和融资。

\begin{warningbox}{AI for Math 的评估不能只看融资}
大额融资和顶级数学家加入是信号，但不是结果。真正结果要看是否能产出可验证证明、加速数学家工作、形成可复用定理库和研究协作模式。
\end{warningbox}

\subsection{本章小结}

AI for Math 的路线从做题走向研究协作。越靠近研究，越需要数学家的直觉、形式化系统和长期反馈，而不是单一 benchmark。

